% =======> Milestone report for the CSE 151B SP26 Math Reasoning Competition.
% Authors: Ketan Mittal, Anirudh Annabathula, Kaii Bijlani.
% Build: pdflatex milestone_report.tex && bibtex milestone_report && pdflatex milestone_report.tex (x2)
% Note: references.bib is referenced but not yet populated; add entries before final submission.
% =================================================

\documentclass{article}

% to avoid loading the natbib package, add option nonatbib:
% \usepackage[nonatbib]{neurips_2024}
% \PassOptionsToPackage{square,comma,sort,numbers}{natbib}
\PassOptionsToPackage{numbers, compress}{natbib}
% \usepackage[numbers]{natbib}

% to compile a preprint version, e.g., for submission to arXiv, add add the
% [preprint] option:
% \usepackage[nonatbib, preprint]{neurips_2024}

% Plain article class is being used (NeurIPS style is commented out above),
% so we load natbib explicitly here to keep \cite/\bibliographystyle{plainnat}
% working.
\usepackage[numbers,compress]{natbib}

\usepackage[utf8]{inputenc} % allow utf-8 input
% \usepackage[T1]{fontenc}    % use 8-bit T1 fonts
% For maths/theorems and such
\usepackage{amsmath}
\usepackage{amssymb}
\usepackage{mathtools}
\usepackage{amsthm}
\usepackage{float}
% Use DeclareMathOperator and/or \newcommand to define abbreviations, e.g.:
\DeclareMathOperator{\R}{\mathbb{R}} % real numbers
% Recommended, but optional, packages for figures and better typesetting:
\usepackage{microtype}
\usepackage{graphicx}
% \usepackage{subfigure} % deprecated; subcaption (below) is the modern replacement
\usepackage{subcaption}
\usepackage{booktabs} % for professional tables (https://www.ctan.org/pkg/booktabs)
% \usepackage{algorithm}  % For algorithms (not used; uncomment if needed)
% \usepackage{algorithmic}
\usepackage{hyperref} % hyperref makes hyperlinks in the resulting PDF.
% \usepackage{wrapfig}  % For wrapping figures with text (not used)
% \usepackage[capitalize,noabbrev]{cleveref}   % For \cref (not in TeXLive basic; using \ref instead)
% \RequirePackage{doi} % not used
% Color references and URLs
\usepackage{xcolor, colortbl}
\definecolor{darkblue}{rgb}{0.19, 0.19, 0.62}
\hypersetup{
  colorlinks = true,
  citecolor  = darkblue,
  filecolor  = darkblue,
  urlcolor   = darkblue,
    % linkcolor  = myred,
}
% Todonotes is useful during development; uncomment if you want \todo notes.
% \usepackage[textsize=tiny]{todonotes}

% \input{_macros}

\title{
CSE 151B Milestone Report --- Team Name
\\
% {\small $<$Your-accessible-and-up-to-date-GitHub-URL$>$} % Note: You can put it in your abstracts or as a footnote in the first page, if you prefer. Note 2: Use \href or \url to link your URL.
}

\author{%
  Ketan Mittal \\
  Department of Data Science \\
  University of California, San Diego \\
  \texttt{k1mittal@ucsd.edu} % TODO: confirm email
  \and
  Anirudh Annabathula \\
  Department of Data Science \\
  University of California, San Diego \\
  \texttt{aannabathula@ucsd.edu} % TODO: confirm email
  \and
  Kaii Bijlani \\
  Department of Data Science \\
  University of California, San Diego \\
  \texttt{kbijlani@ucsd.edu} % TODO: confirm email
}

\begin{document}
\maketitle


\begin{abstract}
This is an in-progress milestone report for our CSE 151B Spring
2026 Math Reasoning Competition entry. The competition gives us 1{,}126
public math problems (375 multiple-choice and 751 free-form) and a
private test set, and asks a language model to answer all of them. So
far we have the inference pipeline working end-to-end: we load
Qwen3-4B-Thinking-2507 with vLLM and INT8 quantization on a single GPU,
generate responses in batches, and score them with the provided
\texttt{judger.py}. The only method change we have made so far is a
new system prompt that tries to match two behaviors of the judger
that were not obvious from the starter code: a $10^{-8}$ relative
numerical tolerance, and a boxed-answer extractor that only
looks at the text after \texttt{</think>} and only keeps the last
contiguous group of boxed expressions. On a 20-question test subset
this prompt change raised accuracy from 50\% (10/20) to 55\%
(11/20), about a 5 percentage point gain. The clearest example is a
temperature-conversion problem where the old prompt produced
decimals like $335.93$, which fail the strict tolerance, and the
new prompt produces exact fractions like $\tfrac{60467}{180}$,
which the judger accepts. We are treating this as preliminary: the
bigger lever we are working on next is GRPO fine-tuning of the base
model on a Claude-filtered subset of the public dataset, where we
use Claude to flag items whose published gold answer looks wrong or
ambiguous and train only on the rest. We did submit our current
pipeline to the Kaggle leaderboard and scored \textbf{0.515} on the
private set, which is roughly in line with the 55\% we saw on the
dev subset and which we treat as a weak starting baseline rather
than a real result.
\end{abstract}


% % % === Introduction
\section{Introduction}
\label{sec:intro}

\paragraph{Problem Definition.}
Each test item is either a multiple-choice question (MCQ) with up to
ten labeled options, or a free-form question that wants one or more
numeric or symbolic values. The model gets the question text (and the
option list when there is one), and has to produce a final answer
wrapped in a boxed expression. The provided \texttt{judger.py} then
parses and scores that answer. For free-form items the judger uses a
sympy-based comparator with a $10^{-8}$ relative tolerance; for MCQ
items it just expects a single capital letter inside the box. So the
task is really a constrained generation problem: the model has to do
multi-step math reasoning \emph{and} put the final answer in a
specific format the judger can parse.

\paragraph{Problem Significance.}
Math reasoning is one of the things people most want LLMs to do
well --- tutoring tools, science assistants, and many engineering
workflows assume the model can both compute the right answer and
report it without ambiguity. The gap between a model that "gets"
a problem internally and one that emits an answer the grading code
will accept is a real, practical robustness problem, and that gap
is exactly what this competition exposes.

\paragraph{Technical Challenge.}
This project is harder than it might look for a few reasons.
(1) \textbf{Data:} the public set covers algebra, calculus,
statistics, probability, linear algebra, geometry, number theory, and
discrete math, and many of the LaTeX strings are slightly corrupted
(e.g., \texttt{frac} or \texttt{int} without backslashes).
(2) \textbf{Grader:} the $10^{-8}$ \emph{relative} tolerance means
that a visibly correct decimal like $0.333$ is graded wrong against
gold $\tfrac{1}{3}$; the boxed-extraction code only keeps the last
contiguous group of boxed expressions in the segment after
\texttt{</think>}. Either of these can quietly cost points if the
prompt does not address them. (3) \textbf{Model:} 4B parameters is
small for math; the model can lose the thread on harder items, and it
likes to restate decimals mid-solution.
(4) \textbf{Engineering:} we have one consumer-grade GPU, so we use
INT8 quantization through BitsAndBytes to fit long thinking traces
(up to 32{,}768 tokens) into the KV cache.
(5) \textbf{Evaluation:} a full pass over 1{,}126 questions is slow,
so we iterate on small subsets and only do full runs occasionally.

\paragraph{Contributions.}
The starter notebook \texttt{starter\_code\_cse151b\_comp.ipynb}
already gives us an end-to-end skeleton: \texttt{uv} environment
setup, vLLM model loading with INT8 quantization, a minimal
MCQ/free-form prompt pair, batched generation, and scoring against
the public set. The main thing it does not do is take the judger's
parsing rules seriously --- the starter prompts only say "put your
answer inside $\backslash$boxed\{\}" and never mention the
$10^{-8}$ tolerance or the contiguous-group extraction rule. So
far our work just addresses that gap; the rest is still in progress.

\begin{itemize}
  \item \textbf{Contribution A (this milestone):} a grader-aware
  system prompt for both free-form and MCQ items. It tells the model
  to keep things in exact symbolic form, to only use a decimal when
  the exact answer is itself a finite decimal (or the problem asks
  for one), and to put multi-part answers in a single
  boxed expression so the judger's contiguous-group rule is
  satisfied. On the 20-question test subset this took overall
  accuracy from 50\% to 55\% and produced a clear qualitative
  win on decimal-vs-exact items (Section \ref{sec:results}).
  \item \textbf{Contribution B (in progress):} scaling the
  evaluation up to the full 1{,}126-item public set, running
  ablations on the prompt's pieces (exact-symbolic clause,
  contiguous-box clause, MCQ recompute-on-mismatch), and trying
  multiple-sample majority voting at the recommended Qwen sampling
  parameters so we can tell prompt effects apart from sampling noise.
  \item \textbf{Contribution C (planned, next major step):} GRPO
  \citep{shao2024deepseekmath} fine-tuning of Qwen3-4B-Thinking-2507
  on a Claude-filtered subset of the public dataset. We follow the
  same recipe the Qwen team uses for their own thinking models in
  the Qwen3 technical report \citep{qwen3technicalreport}, where a
  cold-start SFT stage is followed by GRPO over a filtered set of
  query--verifier pairs. The plan is to ask Claude to read each
  item line by line and flag the ones whose published gold answer
  looks wrong, ambiguous, or underspecified, and then run GRPO on
  the items that survive, using judger correctness as the reward.
  This is meant to attack the kind of errors prompt engineering
  alone cannot fix (real reasoning mistakes on hard items, not
  just formatting drift).
\end{itemize}


% % % === Related Work
\section{Related Work (Optional)}
\label{sec:relatedwork}

\paragraph{LLMs for  mathematical reasoning.}
Open-weight "thinking" models like Qwen3-4B-Thinking-2507
\citep{qwen3technicalreport} are trained to do their reasoning
inside \texttt{<think>...</think>} tokens before answering. The
Qwen team's evaluations on MATH-500-style benchmarks recommend
specific sampling settings ($\text{temp}=0.6$, $\text{top-}p=0.95$,
$\text{top-}k=20$) and a trailing user instruction ("Please reason
step by step, and put your final answer within $\backslash$boxed\{\}").
We use both of those without changes.

\paragraph{GRPO and RL fine-tuning for math.}
Group Relative Policy Optimization (GRPO) was introduced in
DeepSeekMath \citep{shao2024deepseekmath} as a way to fine-tune
reasoning models with a verifier signal, and it showed that using a
deterministic grader as the reward function can noticeably improve
a base model's math accuracy. The same recipe was then used in
Qwen2.5-Math \citep{qwen25math} and again in the Qwen3 technical
report \citep{qwen3technicalreport}, where the team collected
roughly 4{,}000 query--verifier pairs and used GRPO to update the
model after a small cold-start SFT phase. Because we already have a
working judger that returns 0/1 per item and the same backbone
Qwen3 model the technical report describes, that recipe carries
over to our setup almost directly --- which is why we are planning
it as our main next step.

\paragraph{Prompt engineering for structured output.}
A consistent finding in recent LLM-with-grader pipelines is that
spelling out the grader's parsing rules in the system prompt helps
much more than describing the answer's semantics alone. Our
Contribution A is an instance of this idea, specialized to the
$10^{-8}$ relative tolerance and the contiguous-boxed-group rule
of the provided judger.


% % % === Methods
\section{Methods}
\label{sec:methods}

\paragraph{Problem Setting.}
We treat each item as a tuple $(q, O, y)$ where $q$ is the
question text, $O$ is the option list (empty for free-form items),
and $y$ is the gold answer (a capital letter for MCQ; a list of one
or more numeric/symbolic strings for free-form). The model defines a
conditional distribution $p_\theta(r \mid s(q, O), u(q, O))$ over
response strings $r$, where $s$ and $u$ are the system and user
messages built by \texttt{build\_prompt}. The judger
$\mathcal{J}(r, y, O) \in \{0, 1\}$ is the function in
\texttt{judger.py}: it pulls the last contiguous group of
boxed expressions from $r$ (only looking at the text after
\texttt{</think>}), normalizes the LaTeX, and compares to $y$
using sympy with a $10^{-8}$ relative tolerance on numerical
values. We want to maximize the expected accuracy
$\mathbb{E}_{(q, O, y) \sim \mathcal{D}}\mathbb{E}_{r \sim p_\theta}[\mathcal{J}(r, y, O)]$.
At this milestone we do \emph{not} change $\theta$ (no fine-tuning
yet); we only change $s$. Fine-tuning is planned for the next phase
(see Contribution C).

The dataset is \texttt{data/public.jsonl} (1{,}126 items: 375 MCQ,
751 free-form). For fast iteration we use the first 20 items as a
dev subset. The model sees the raw question and option strings, and
we keep its full response as-is. We do not pre- or post-process the
strings beyond what the judger already does at scoring time.

\paragraph{Idea Summary.}
The starter prompt is short: "solve step by step" and "put your
final answer inside $\backslash$boxed\{\}." We noticed it does
not mention two things the judger actually cares about:
(i) \textbf{decimal drift}, where the model gets $\tfrac{565}{9}$
internally but writes "$\approx 62.78$" inside the box and fails
the $10^{-8}$ tolerance against the exact fraction; and
(ii) \textbf{boxed fragmentation}, where the model writes a final
non-contiguous box after restating an earlier answer and the judger
silently drops everything except the last group. Our new prompt
addresses both. For free-form items we explicitly tell the model to
prefer exact symbolic forms (e.g., $\boxed{\tfrac{1}{3}}$ over
$\boxed{0.333}$, $\boxed{\sqrt{2}}$ over $\boxed{1.414}$),
list the carve-outs (decimals are fine when the exact answer is
already a finite decimal, or when the problem explicitly asks for
one), and ban $\backslash\text{approx}$ inside the box. For
multi-part answers we require a single boxed expression with
comma-separated values in the order asked, which is the easy way
to never trip the contiguous-group rule. For MCQ items we tell the
model to compute the exact symbolic answer first and only convert to
a decimal at the end to match against the options, and to recompute
from scratch (rather than picking the closest-looking option) when
nothing matches.

There is a real risk that this prompt is too prescriptive and ends
up distracting the model from the actual math. We try to mitigate
that by keeping the reasoning instructions short and putting all the
hard format rules in their own "Final-answer format" section so the
model can read past them once it knows the rules. If the prompt does
not generalize past the dev subset we have a fallback: take the
model's response and re-prompt it to re-cast just the final box in
exact symbolic form, which keeps the original reasoning intact and
only patches the formatting.

After we have run the prompt-engineering ablations in earnest, the
next planned step is GRPO fine-tuning. The setup we have in mind is
to first run Claude over the public dataset to filter out items
whose gold answer is wrong, ambiguous, or underspecified, then run
GRPO on the surviving items with the judger as the reward function
($+1$ for correct, $0$ otherwise). The Claude-based filter is
important because training a small model with RL against noisy gold
answers tends to make things worse, not better, so we want to make
sure the supervision signal is clean before we update any weights.


% % % === Experiments
\section{Experiments}
\label{sec:experiments}

\subsection{Baselines}
\label{sec:baselines}
We compare two prompts. Everything else --- model, quantization,
sampling parameters, dataset subset --- is the same for both runs:
\begin{itemize}
    \item \textbf{Starter prompt} (baseline, from
    \texttt{starter\_code\_cse151b\_comp.ipynb}): a two-sentence
    system message that asks the model to solve step by step and put
    the final answer "inside $\backslash$boxed\{\}" with
    comma-separated values for multi-answer items. It says nothing
    about decimal vs.\ exact form or about the contiguous-group rule,
    and the MCQ version just says "output ONLY the letter."
    \item \textbf{Refined prompt (ours)}: the
    \texttt{SYSTEM\_PROMPT\_MATH} and \texttt{SYSTEM\_PROMPT\_MCQ}
    we describe in Section \ref{sec:methods}, plus the Qwen team's
    recommended user-side suffix "Please reason step by step, and
    put your final answer within $\backslash$boxed\{\}."
\end{itemize}
We do not include simpler models like a multi-layer perceptron or
linear regression as baselines, because they cannot really do
free-form math reasoning over natural-language problem statements
in any meaningful way. Comparing prompt variants of the same backbone
model is what made sense to us at this stage.

\subsection{Evaluation}
\label{sec:evaluation}
We score with the provided judger: \texttt{Judger.auto\_judge} for
free-form items, and an exact-letter match (after extracting the
boxed letter) for MCQ items. This is the same scoring the starter
notebook uses. For this milestone we report accuracy on the first
20 items of \texttt{public.jsonl}, split by question type. We are
saving the full public set and any Kaggle submissions for the final
report. We also do qualitative case studies on items whose
correctness flipped between the two prompts, since those are the
most informative about what the prompt change is and is not buying
us.

\subsection{Implementation details}
\label{sec:implementation_details}
\begin{itemize}
    \item \textbf{Hardware.} A single CUDA GPU
    (\texttt{CUDA\_VISIBLE\_DEVICES=0}). Generating responses for
    the 20-question subset takes about 2--3 minutes after the model
    is loaded; we estimate a full-public-set run (1{,}126 questions)
    at roughly 90--120 minutes depending on how long the responses
    end up being. We are not training yet, so there is no per-epoch
    time to report.
    \item \textbf{Model and quantization.} \texttt{Qwen/Qwen3-4B-Thinking-2507}
    loaded through vLLM with \texttt{quantization="bitsandbytes"} and
    \texttt{load\_format="bitsandbytes"} (INT8 weights),
    \texttt{gpu\_memory\_utilization=0.85},
    \texttt{max\_model\_len=8192},
    \texttt{max\_num\_seqs=256},
    \texttt{max\_num\_batched\_tokens=32768},
    \texttt{enable\_prefix\_caching=False},
    \texttt{trust\_remote\_code=True}.
    \item \textbf{Sampling.} We have not deliberately tuned sampling
    at this milestone. Our active inference cell sets
    \texttt{temperature=0.6}, \texttt{top\_p=0.95}, \texttt{top\_k=20},
    \texttt{repetition\_penalty=1.0}, and \texttt{max\_tokens=32{,}768},
    and otherwise leaves things at their defaults. These values
    happen to line up with what the Qwen team recommends for their
    thinking models, but we have not yet run a controlled comparison
    against alternatives --- intentionally validating the
    recommended settings (and varying them to see what actually
    matters) is on our list for the next phase.
    \item \textbf{Keeping the comparison fair.} Both prompt variants
    use the same sampling parameters, the same chat template (the
    tokenizer's default \texttt{apply\_chat\_template} with
    \texttt{add\_generation\_prompt=True}), and the same scoring
    code. The only thing that changes between the two runs is the
    contents of \texttt{SYSTEM\_PROMPT\_MATH} and
    \texttt{SYSTEM\_PROMPT\_MCQ}.
    \item \textbf{Caveat.} We do not fix a random seed for vLLM at
    this milestone, so per-question outcomes have some sampling
    noise. Section \ref{sec:results} talks about how this shows up.
\end{itemize}

\subsection{Results}
\label{sec:results}

\paragraph{Result 1: the new prompt raises accuracy by $+5$ percentage points on the dev subset.}
Table \ref{tab:dev-subset} shows the head-to-head on the 20-question
test subset. Overall accuracy goes from 50\% (10/20) to 55\%
(11/20). The MCQ split is unchanged at 4/9; the free-form split
improves from 6/11 to 7/11.

\begin{table}[h]
  \caption{Accuracy on the 20-question test subset (first 20 items
  of \texttt{public.jsonl}). Both rows use Qwen3-4B-Thinking-2507
  (INT8 via vLLM) with the same sampling parameters; only the system
  prompt is different.}
  \label{tab:dev-subset}
  \centering
  \begin{tabular}{lccc}
    \toprule
    Prompt          & MCQ ($n=9$) & Free-form ($n=11$) & Overall ($n=20$) \\
    \midrule
    Starter         & 4 / 9         & 6 / 11               & 10 / 20 \;(50.0\%) \\
    Refined (ours)  & 4 / 9         & 7 / 11               & \textbf{11 / 20 \;(55.0\%)} \\
    \bottomrule
  \end{tabular}
\end{table}

\paragraph{Result 2: a clear qualitative example of decimal-vs-exact.}
The clearest example of what the new prompt is doing is item
\texttt{id=5}, a free-form temperature-conversion problem with gold
$=[62.7777777777778,335.927777777778,604.67]$. With the starter
prompt the model writes
$\boxed{62.78,335.93,604.67}$; the middle value differs from
gold by about $10^{-4}$ relative, which is way outside the
$10^{-8}$ tolerance, so the whole item is marked wrong. With the
refined prompt the model writes
$\boxed{\tfrac{565}{9},\tfrac{60467}{180},\tfrac{60467}{100}}$;
each fraction equals the gold decimal exactly under sympy, so the
item is marked correct. This is exactly the failure mode the new
prompt was designed to fix, and seeing it work on a real item is
the main reason we believe the prompt change is helping.

\paragraph{Result 3 (negative finding): the gain is small enough that sampling noise still matters.}
Three of the 20 items flipped between the two runs: \texttt{id=5}
(free-form, F$\rightarrow$T) and \texttt{id=19} (MCQ,
F$\rightarrow$T) are wins, but \texttt{id=13} (MCQ, T$\rightarrow$F)
is a regression. Item 13 is a triple-integral MCQ; the refined-prompt
run loses the thread mid-reasoning rather than producing a malformed
final answer, so we do not think the prompt is actually making the
model worse there --- it is more likely just sampling variance, since
temperature is 0.6 and we did not seed. The honest read of the
$+1$ net gain on $n=20$ is that it is a noisy point estimate, and
we plan to tighten it in the final report by (a) running on the
full public set and (b) averaging over multiple sampling seeds.


% % % === Discussion
\section{Discussion}
\label{sec:discussion}
At this point in the project we have an inference pipeline that runs
end-to-end (model loading, batched generation, judger-based scoring)
and one method change --- the new system prompt --- that removes a
specific kind of avoidable loss (decimal drift on free-form items).
Most of the work is still ahead of us, and the rest of this section
is about what we still need to do to make it better.

\paragraph{Achievements.}
We got the starter pipeline running locally on our own hardware, read
through \texttt{judger.py} carefully enough to understand the two
behaviors that surprised us (the $10^{-8}$ relative tolerance and
the post-\texttt{</think>} contiguous-boxed-group extraction), and
used that understanding to design a new system prompt that explicitly
respects both. On the 20-question test subset the refined prompt
gave a $+5$ percentage point gain (50\% $\to$ 55\%) under
matched sampling parameters, with the win coming from the free-form
split ($+1$ of 11). We also submitted the refined pipeline's
private-set predictions (943 items) to Kaggle and scored
\textbf{0.515} on the private leaderboard. This is honestly not a
strong result: it lines up with the 55\% we saw on the dev subset
and confirms that prompt engineering alone is not enough to get us
where we want to be. We are reporting it as our current baseline
rather than as a finished result --- the rest of this section is
about why the score is where it is and what we are doing next to
improve it.

\paragraph{Bottlenecks and Mitigation.}
The biggest issue right now is statistical: 20 questions is just not
enough to tell prompt effects apart from sampling noise (see
Result 3). The plan is to run on the full 1{,}126-item public set and
average over multiple sampling seeds for the final report. The
second issue is wall-clock time --- a full-public-set pass takes
roughly 90--120 minutes on our hardware. We are planning to do a
single full sweep and then iterate on cheaper post-hoc rewrites
rather than re-running inference each time. The third is the
LaTeX corruption in some of the public-set questions
(\texttt{frac} without a backslash, etc.); the judger's
\texttt{normalize\_answer} handles most of it, but if we see a
recurring error class we will add a small pre-processor to the
question text.

\paragraph{Plans Forward.}
Between now and the final report we plan to:
\begin{itemize}
    \item Run a full 1{,}126-item public-set evaluation under both
    prompts so the comparison is on real numbers, not 20.
    \item Ablate the new prompt's individual pieces (exact-symbolic
    clause, contiguous-box clause, MCQ recompute-on-mismatch clause)
    to figure out which parts are actually doing the work.
    \item Deliberately validate the Qwen team's recommended
    thinking-mode sampling settings. We currently have those values
    in our code but have not run a controlled comparison; we plan to
    sweep \texttt{temperature}, \texttt{top\_p}, \texttt{top\_k}, and
    \texttt{min\_p} on a fixed subset and check whether the
    recommended values actually help in our pipeline.
    \item Try multi-sample majority voting at $k \in \{3, 5\}$,
    applying the judger's boxed extractor to each sample
    independently. This should reduce the sampling noise we saw in
    Result 3.
    \item Build a small post-hoc rewriter that takes the model's
    response and re-emits just the final box in clean exact symbolic
    form. That way reasoning errors and formatting errors stop being
    coupled.
    \item \textbf{(Main next step) Run GRPO fine-tuning on a
    Claude-filtered subset of the public dataset.} We use Claude as
    a verifier: for each public-set item, Claude reads the question,
    the published gold answer(s), and the listed options, and tells
    us whether the gold answer is plausibly correct, ambiguous, or
    flat-out wrong. We drop everything Claude flags as wrong or
    ambiguous and keep only the high-confidence items. Then we run
    GRPO on Qwen3-4B-Thinking-2507 with that filtered set, using the
    judger's 0/1 correctness signal as the reward. The hope is that
    this catches the errors prompt engineering cannot fix --- real
    reasoning mistakes on hard problems, not just bad formatting.
\end{itemize}
This is broadly consistent with the original proposal (prompt
engineering as the first step, fine-tuning as the next). The two
deviations from the proposal are (a) the multi-sample voting plan,
which we added once we saw the sampling noise in Result 3, and
(b) the Claude-based dataset filter, which we added once we noticed
some of the public-set gold answers look suspect.

\paragraph{Team Responsibilities.}
\begin{itemize}
    \item \textbf{Ketan Mittal} --- grader analysis and prompt
    design; ran the starter-vs-refined comparison; produced the
    private-set prediction file.
    % TODO: confirm and expand
    \item \textbf{Anirudh Annabathula} --- inference infrastructure
    (vLLM, INT8 quantization, batched generation); will own the
    GRPO training setup.
    % TODO: confirm and expand
    \item \textbf{Kaii Bijlani} --- dataset analysis and qualitative
    case-study pipeline; will own the Claude-based dataset filter.
    % TODO: confirm and expand
\end{itemize}

\paragraph{Timelines.}
\begin{itemize}
    \item \textbf{Week 2:} Form team; reproduce the starter notebook
    locally.
    \item \textbf{Week 3:} Read and document \texttt{judger.py};
    identify the $10^{-8}$ relative-tolerance and
    contiguous-boxed-group behaviors.
    \item \textbf{Week 4:} Draft and check the grader-aware system
    prompt on the 20-question test subset; produce the results in
    Table \ref{tab:dev-subset}; generate the first private-set
    prediction file.
    \item \textbf{Week 5 (current):} Submit milestone report.
    \item \textbf{Week 6:} Run the full 1{,}126-item public-set
    evaluation under both prompts; record per-question correctness
    so we can ablate later.
    \item \textbf{Week 7:} Prompt-component ablations; multi-sample
    majority-voting experiments.
    \item \textbf{Week 8:} Run Claude as a verifier over the public
    set and build the filtered training subset; stand up the GRPO
    training loop.
    \item \textbf{Week 9:} GRPO fine-tuning runs; lock final
    results; resubmit to Kaggle and try to beat the current 0.515
    baseline.
    \item \textbf{Week 10:} Write up and submit the final report.
\end{itemize}


\section*{LLM Usage}
We used an LLM assistant (Anthropic Claude) for two things: (i) help
with wording, mostly tightening up phrasing in this report and in our
system prompts, and (ii) help making sense of hard-to-read error
messages in LaTeX and Python while we were writing the report and
building the pipeline. All the technical decisions, code,
experimental runs, and numbers in this report are ours, and we take
full responsibility for the contents.


% % % === ACKNOWLEDGMENTS AND REFERENCES START
% \section*{Acknowledgements}
% This work was supported $\dots$.
\bibliography{references.bib}
\bibliographystyle{plainnat}  % Can also use abbrvnat, unsrt, etc.
%
% % % === APPENDIX (optional)
% \newpage
% \appendix
% \section*{Appendix}
% \section{Appendix topic 1}
%
\end{document}
